\documentclass[11pt]{letter}

\usepackage[margin=0.75in]{geometry}
\usepackage{setspace}
\usepackage{hyperref}
\usepackage{xcolor}
\usepackage{enumitem}
\usepackage{parskip}

\hypersetup{
    colorlinks=true,
    urlcolor=black,
    linkcolor=black
}

\setlength{\parindent}{0pt}
\setstretch{1.05}
\date{}
\begin{document}

\begin{letter}{}

\opening{Dear Hiring Committee,}

I saw the job post on \textbf{Planet's careers site} and thought I'd be a good fit for it.

I am a Master's student in Computer Science at the University of Cincinnati, graduating December 2026, looking for a role as a \textbf{Software Engineer, Enterprise AI} at \textbf{Planet}. With hands on experience building backend automation, cloud platforms on AWS and GCP, and AI integrated applications using RAG pipelines and agent tool calling, I am excited about the opportunity to contribute to Planet's Corporate Engineering team and its work scaling enterprise AI capabilities across the organization.

\textbf{Here are reasons that make me a great fit for the role:}

\begin{itemize}[leftmargin=*, itemsep=0.45em]
    \item \textbf{Enterprise AI application development:} I architected an evidence grounded RAG pipeline over Google Gemini 2.5 flash for PolicyPilot, forcing citation backed answers to eliminate hallucination in high stakes healthcare eligibility decisions, and designed a 17 tool MCP server exposing structured LLM tool calling consumable by external AI agents, similar in spirit to cataloguing and governing agents across a company.

    \item \textbf{AI orchestration and agentic workflows:} In my graduate research, I built a provider agnostic LLM abstraction layer routing across Gemini, GPT-4o, Llama, and Qwen, and designed a two stage memory augmented agent pipeline for Android bug reproduction, achieving a 25.7x latency reduction over a published VLM baseline while preserving task completion fidelity.

    \item \textbf{Hands on GCP and AWS experience in regulated environments:} At Siemens Industry Software, I automated Teamcenter X infrastructure workflows with Python and Boto3 across AWS environments, cutting deployment time by 40 percent, and managed HashiCorp Vault secrets and IAM access policies to strengthen data access controls, work that closely mirrors the compliance and secret management rigor Planet's AI integrations require.

    \item \textbf{Python, vector databases, and AI orchestration platforms:} Across PolicyPilot and Spendora, I have built RAG pipelines with LangChain, integrated vector based retrieval for grounded LLM responses, and worked with Python, SQL, and REST APIs to serve real time AI driven insights.

    \item \textbf{Cross team collaboration and evangelism:} As Quality Coach and SonarQube Specialist at Siemens, I increased automated code coverage from 35 percent to 78 percent and mentored 5 software engineers, experience that reflects the cross functional partnership and best practice evangelism this role calls for.
\end{itemize}

I am available for full-time work after December 2026. As a Master's student at the University of Cincinnati and someone who values professionalism and growth, I am enthusiastic about contributing to Planet's mission of using space to help life on Earth.

Please find my CV attached below. I look forward to hearing from you. Planet's work turning satellite data into tools that solve real world problems is exactly the kind of mission driven engineering I want to be part of.

{\setlength{\parskip}{0pt}
Regards,\\
Tanmay Bhuskute,\\
bhuskutv@mail.uc.edu $\mid$ tanmay.v.bhuskute@gmail.com $\mid$ +1 7328583027,\\
3424 Brookline Ave, Flat 19, Cincinnati, OH 45220.
}

\end{letter}

\end{document}
